% PDF-to-TeX transcription by the assistant; not the original downloaded file.
\begin{theorem}[Section 2.4.6]
The effective resistance $R_{\mathrm{eff}}$ from $0$ to $\infty$ is less than or equal to the energy dissipation of any unit flow from $0$ to infinity.
\end{theorem}
\begin{proof}
Assume that we have a unit flow $j$ from $0$ to infinity with energy dissipation
\[
E=\frac12\sum_{x,y}j_{xy}^2R_{xy}.
\]
We claim that $R_{\mathrm{eff}}\le E$. Restricting $j_{xy}$ to the edges of the finite graph $G(r)$, we have a unit flow from $0$ to $\partial G(r)$ in $G(r)$. Let $i^{(r)}$ be the unit current flow in $G(r)$ from $0$ to $\partial G(r)$. By the results of Section 1.3.5,
\[
R_{\mathrm{eff}}^{(r)}=\frac12\sum_{G(r)}(i^{(r)}_{xy})^2R_{xy}
\le\frac12\sum_{G(r)}j_{xy}^2R_{xy}
\le\frac12\sum_{x,y}j_{xy}^2R_{xy}=E,
\]
where $\sum_{G(r)}$ indicates the sum over all pairs $x,y$ such that $xy$ is an edge of $G(r)$.
\end{proof}
\paragraph{A proof, using flows, that simple random walk in three dimensions is transient (Section 2.4.7).}
We now apply this form of the cutting method to give another proof that simple random walk on the threedimensional lattice is transient. All we need is a flow to infinity with finite dissipation.
\begin{proof}[Three-dimensional construction and energy estimate]
For three dimensions, the uniform flow is defined as follows: Out of $(x,y,z)$ with $x+y+z=n$ we have the flow indicated in Figure 71. The total flow out of $(x,y,z)$ is then
\[
\frac{2(x+1)}{(n+3)(n+2)(n+1)}+
\frac{2(y+1)}{(n+3)(n+2)(n+1)}+
\frac{2(z+1)}{(n+3)(n+2)(n+1)}
=\frac2{(n+2)(n+1)}.
\]
\begin{center}
\begin{tabular}{c|c}
\multicolumn{2}{c}{Figure 71: edge labels transcribed from the original diagram}\\
Outgoing direction & Flow \\\hline
positive $x$ & $\displaystyle\frac{2(x+1)}{(n+3)(n+2)(n+1)}$\\[7pt]
positive $y$ & $\displaystyle\frac{2(y+1)}{(n+3)(n+2)(n+1)}$\\[7pt]
positive $z$ & $\displaystyle\frac{2(z+1)}{(n+3)(n+2)(n+1)}$
\end{tabular}
\end{center}
The flow into $(x,y,z)$ comes from the points $(x-1,y,z)$, $(x,y-1,z)$, $(x,y,z-1)$ and, hence, the total flow into $(x,y,z)$ is
\[
\frac{2x}{(n+2)(n+1)n}+
\frac{2y}{(n+2)(n+1)n}+
\frac{2z}{(n+2)(n+1)n}
=\frac2{(n+2)(n+1)}.
\]
Thus the net flow for $(x,y,z)$ is $0$. The flow out of $0$ is $(1/3)+(1/3)+(1/3)=1$ We have now to check finiteness of energy dissipation. The flows coming out of the edges at the $n$th level are all $\le2/(n+1)^2$. There are $(n+1)(n+2)/2$ points a distance $n$ from $0$, and thus there are $(3/2)(n+1)(n+2)\le3(n+1)^2$ edges coming out of the $n$th level. Thus the energy dissipation $E$ has
\[
E\le\sum_n3(n+1)^2\left(\frac2{(n+1)^2}\right)^2
=12\sum_n\frac1{(n+1)^2}<\infty,
\]
and the random walk is transient.
\end{proof}
